\section{The return order when \texorpdfstring{$n=3m+1$}{n=3m+1}}\label{app:cycles}

We calculate the order of the map $M_1$ in \eqref{eq:efficient-map} for $m\ge3$.
\begin{table}[htbp]
\centering
\renewcommand{\arraystretch}{1.3}
\begin{tabular}{c l}
\toprule
$m\bmod6$ & Cycle lengths defining $R_1$\\
\midrule
$0,2$ & $3m+1$\\
$1$ & $(2m+1)/3,\ (7m+2)/3$\\
$3$ & $(7m+3)/6,\ (11m+3)/6$\\
$4$ & $(2m+1)/3,\ (7m+2)/6,\ (7m+2)/6$\\
$5$ & $(7m+1)/6,\ (11m+5)/6$\\
\bottomrule
\end{tabular}
\caption{The nonzero cycles of the return map when $n=3m+1$.
The order $R_1$ is their least common multiple.}
\label{tab:cycles}
\end{table}

The nonzero cycles of $g$ are
\[
 (3,6,\ldots,3m),\qquad
 (1,5,8,\ldots,3m-1,2,4,7,\ldots,3m+1).
\]
The first has length $m$ and the second length $2m+1$. To evaluate
$g^{m-1}$ on the second, index its letters by
\[
 B_0=1,\quad B_j=3j+2\ (1\le j<m),\quad B_m=2,\quad
 B_{m+j}=3j+1\ (1\le j\le m).
\]
The map $g$ adds one to these indices modulo $2m+1$. Also
\[
 (\rho g^{-1})(1)=2,\quad
 (\rho g^{-1})(2)=3m,\quad
 (\rho g^{-1})(3)=3m+1,\quad
 (\rho g^{-1})(4)=3,\quad
 (\rho g^{-1})(5)=1,
\]
and $(\rho g^{-1})(x)=x-2$ for $6\le x\le3m+1$.
Applying $g^{m-1}$ and then $\rho$ gives
\[
\begin{array}{c|rrrrrrrrr}
x&1&2&3&4&5&6&7&8&9\\\hline
M_1(x)&3m-1&3m-2&3m-3&3m+1&3m&2&3&4&1,
\end{array}
\]
and, for $10\le x\le3m+1$,
\begin{equation}\label{eq:descent}
 M_1(x)=\begin{cases}
 x-6,&x\equiv0\pmod3,\\
 x-5,&x\equiv1\pmod3,\\
 x-4,&x\equiv2\pmod3.
 \end{cases}
\end{equation}
For example, when $x=3i\ge6$, its image under $\rho g^{-1}$ is
$B_{m+i-1}$. Adding $m-1$ to this index gives $2m$ at $i=2$, zero
at $i=3$, and $i-3$ for $i\ge4$. These are the exceptional values
$M_1(6)=2$, $M_1(9)=1$, and the first line of \eqref{eq:descent}.
The other two residues follow by the same index substitution.

Every orbit outside $\{1,2,3,4,5\}$ decreases until it reaches that
set. Thus every cycle contains one of these five \emph{anchors}.
The descent times are as follows; the first jump out of an anchor is
not included.
\[
\renewcommand{\arraystretch}{1.15}
\begin{array}{c|c|c}
\text{Starting value}&\text{Anchor reached}&\text{Steps}\\\hline
6j\ (j\ge1)&2&j\\
6j+3\ (j\ge1)&1&j\\
9j+1\ (j\ge1)&5&2j-1\\
9j+4\ (j\ge0)&4&2j\\
9j+7\ (j\ge0)&3&2j+1\\
9j+2\ (j\ge1)&3&2j\\
9j+5\ (j\ge0)&5&2j\\
9j+8\ (j\ge0)&4&2j+1
\end{array}
\]
For multiples of three, this is repeated subtraction of six, ending
at $6\mapsto2$ or $9\mapsto1$. In the other residues, two steps
subtract nine. The low values $4,5,7,8,10$ give the stated endpoints.

Write $m=6h+r$. Applying the descent times to the five initial jumps
gives the following first-return data, in anchor order $1,2,3,4,5$.
\[
\renewcommand{\arraystretch}{1.2}
\begin{array}{c|c|c}
r&\text{Destinations}&\text{Arc lengths}\\\hline
0&(4,3,1,5,2)&(4h,4h,3h,4h,3h+1)\\
1&(3,5,2,4,1)&(4h+1,4h,3h+1,4h+1,3h+1)\\
2&(5,4,1,3,2)&(4h+1,4h+1,3h+1,4h+2,3h+2)\\
3&(4,3,2,5,1)&(4h+2,4h+2,3h+2,4h+2,3h+2)\\
4&(3,5,1,4,2)&(4h+3,4h+2,3h+2,4h+3,3h+3)\\
5&(5,4,2,3,1)&(4h+3,4h+3,3h+3,4h+4,3h+3).
\end{array}
\]
Summing the arc lengths around each anchor cycle gives exactly
Table~\ref{tab:cycles}. There are no further cycles, since every
nonanchor reaches an anchor.

For a single cycle, the order is $n\le n^2/4$. For two cycles of
lengths $u,v$, it is at most $uv\le n^2/4$. The three-cycle case has
lengths $u,v,v$, so its order is at most $uv\le n^2/8$.
This proves the quadratic bound used in the main theorem.
